\chapter{NKMedia : lire et écrire de la vidéo}

Ce chapitre est le plus court des cinq, et c'est délibéré.

\texttt{NKMedia} est de très loin le plus gros module du moteur. Il contient des
décodeurs H.264, HEVC, VP8, VP9, AV1, MPEG-2, Theora, AAC, Opus (avec ses deux
moitiés CELT et SILK), AMR — tous écrits à la main, sans \texttt{ffmpeg}. Le
seul décodeur AV1 pèse 286~Ko de source. Un chapitre pour débutants qui
essaierait d'en faire le tour serait illisible et faux.

Nous nous limiterons donc à \textbf{quatre classes de façade}, plus une
cinquième en bonus. Tout le reste est de la plomberie interne, et vous n'avez
aucune raison d'y toucher.

\begin{nkretenir}[Le périmètre de ce chapitre]
\begin{itemize}
  \item \texttt{NkMediaProbe} — qu'y a-t-il dans ce fichier ?
  \item \texttt{NkVideoWriter} — écrire une vidéo, en MJPEG ;
  \item \texttt{NkVideoReader} — la lire image par image ;
  \item \texttt{NkVideoRecorder} — enregistrer ce que l'application affiche ;
  \item \texttt{NkImageSequenceWriter} — écrire une séquence d'images.
\end{itemize}
Cinq classes, aucune ligne de codec. Si vous maîtrisez ces cinq-là, vous savez
faire tout ce dont une application a besoin.
\end{nkretenir}

\section{Le module, et une inversion inhabituelle}

\begin{nkcode}{Kernel/Runtime/NKMedia/NKMedia.jenga:17-22}
    nkentseudependson(
        ["NKCore", "NKPlatform", "NKMemory", "NKContainers", "NKMath", "NKLogger",
         "NKStream", "NKFileSystem", "NKImage", "NKThreading", "NKTime"],
        selfexport="NKMedia",
        extra_includes=["src"],
    )
\end{nkcode}

\texttt{NKImage} est dans la liste, et pour une raison élégante : \textbf{le
codec MJPEG n'est rien d'autre que le codec JPEG de NKImage appliqué image par
image}. Les séquences d'images passent, elles aussi, par les codecs PNG, BMP,
TGA et QOI de NKImage. Le chapitre~5 n'était donc pas un préalable arbitraire.

Il n'y a pas d'en-tête parapluie : on inclut le sous-en-tête précis dont on a
besoin. Tout vit dans le \emph{namespace} \texttt{nkentseu::media}.

\begin{nkcode}{Kernel/Runtime/NKMedia/ — arborescence (abrégée)}
NKMedia/
  NKMedia.jenga · ROADMAP.md (1529 l. — la source de verite sur l'etat reel)
  src/NKMedia/
    NkMediaProbe.{h,cpp}        <- detection conteneur + pistes
    NkMediaDemux.{h,cpp}        <- extraction des paquets audio
    Video/
      NkVideoTypes.h            <- NkVideoInputFormat
      NkVideoReader.{h,cpp}     <- LECTURE : conteneur -> RGBA8
      NkVideoWriter.{h,cpp}     <- ECRITURE simple
      NkVideoRecorder.{h,cpp}   <- CAPTURE A/V threadee
      NkVideoConverter.{h,cpp}  <- sequence/AVI -> video
      NkImageSequenceWriter.{h,cpp}
      Containers/ NkAviWriter · NkMovWriter · NkMp4H264Writer · NkWebmWriter
    Audio/Containers/NkWavWriter.{h,cpp}
    Codecs/ Video/{H264,HEVC,VP8,VP9,AV1,Mpeg1,Mpeg2,Theora} · Aac · Opus · AMR
\end{nkcode}

\subsection{Ici, les commentaires \emph{sous-estiment} le module}

Nous avons pris l'habitude, depuis le chapitre~5, de nous méfier des
commentaires qui promettent trop. NKMedia présente le problème inverse, et il
faut le savoir pour ne pas renoncer à des fonctionnalités qui existent.

\begin{nkcode}{Kernel/Runtime/NKMedia/src/NKMedia/Video/NkVideoReader.h:8-11}
//   - Conteneur MOV/MP4 (ISOBMFF) : piste vidéo -> MJPEG (mjpa/jpeg). H264 = à venir.
// H264/AVC (MP4 courant) exige un décodeur complet (gros chantier séparé) : la
// lecture d'une piste avc1 renvoie une erreur claire tant que le décodeur n'existe pas.
\end{nkcode}

C'est faux. Le fichier \nkfile{Video/NkVideoReader.cpp} affecte réellement
\texttt{info.codec} à \texttt{"h264"} (lignes 1592, 1828, 2002), \texttt{"hevc"}
(955, 1563), \texttt{"vp8"} (1572, 1650), \texttt{"vp9"} (1578, 1635),
\texttt{"av1"} (1586, 1644), \texttt{"mpeg2"} (1869), \texttt{"theora"} (1944),
en plus de \texttt{"mjpeg"}, \texttt{"rawrgb"} et \texttt{"image"}. Les
décodeurs sont branchés.

De même, \nkfile{Video/NkVideoWriter.h:5-6} annonce « AVI pour l'instant ;
MP4/WebM à venir » alors que les écrivains MOV/MP4 et WebM sont livrés.

\begin{nkretenir}[La source de vérité de ce module]
Pour NKMedia, la vérité n'est ni dans les en-têtes ni dans les
\texttt{.jenga} : elle est dans \nkfile{Kernel/Runtime/NKMedia/ROADMAP.md}, un
document de 1\,529 lignes tenu à jour. Prenez le réflexe de l'ouvrir avant de
conclure qu'une fonctionnalité manque.
\end{nkretenir}

Voici l'essentiel de cette feuille de route, au moment de la rédaction :

\begin{center}
\begin{tabular}{@{}ll@{}}
\toprule
\textbf{Brique} & \textbf{Statut} \\
\midrule
Probe / démux conteneurs & fini — MP4, WebM, WAV, OGG, MP3, FLAC \\
Extraction de paquets    & fini — MP4 (\texttt{stbl} et fMP4), WebM \\
Opus CELT / SILK         & fini (SILK exact au bit face à libopus) \\
Dispatcher Opus          & \emph{partiel} — mode hybride non fait \\
Décodeur AAC-LC          & fini, exact au bit \\
Muxers (écriture)        & fini — AVI, MOV/MP4, WAV, WebM \\
Vidéo (décodage)         & fini — H.264 Main+High, VP8, VP9, HEVC \\
\textbf{Vidéo (encodage)} & \textbf{en cours} — RAW/MJPEG/MPEG-1 + H.264 baseline \\
AV1 / MPEG-2 / Theora    & \textbf{restent à faire} \\
\bottomrule
\end{tabular}
\end{center}

\begin{nkpiege}[Ce sur quoi ce cours ne s'engage pas]
\begin{itemize}
  \item \textbf{L'encodage H.264} est marqué « en cours » et limité au profil
        \emph{baseline}. Il fonctionne — le \texttt{NkVideoRecorder} l'utilise
        par défaut — mais ce n'est pas un chemin sur lequel bâtir un exercice de
        débutant. \textbf{Restez en MJPEG} ;
  \item \textbf{AV1, MPEG-2 et Theora} : le code de décodage existe, la feuille
        de route ne les déclare pas terminés ;
  \item \textbf{Opus en mode hybride} et en stéréo peut être « refusé
        proprement » : un fichier \texttt{.opus} courant n'est pas garanti ;
  \item \textbf{HEVC} refuse proprement les tuiles, le PCM, les
        échantillonnages 4:2:2 et 4:4:4 ;
  \item \textbf{VP8} refuse la segmentation par macrobloc et les partitions
        multiples.
\end{itemize}
« Refusé proprement » est une bonne nouvelle : cela signifie une erreur claire
plutôt qu'une image corrompue.
\end{nkpiege}

\section{\texttt{NkMediaProbe} : qu'y a-t-il dans ce fichier ?}

C'est la classe par laquelle commencer, parce qu'elle ne décode rien : elle lit
les en-têtes et vous dit ce que contient le fichier.

\begin{nkcode}{Kernel/Runtime/NKMedia/src/NKMedia/NkMediaProbe.h:21-70 (abrégé)}
enum class NkMediaContainer { NK_UNKNOWN, NK_MP4, NK_WEBM, NK_WAV, NK_OGG, NK_MP3, NK_FLAC };
enum class NkMediaTrackType { NK_UNKNOWN, NK_AUDIO, NK_VIDEO };

struct NkMediaTrack {
        NkMediaTrackType type = NkMediaTrackType::NK_UNKNOWN;
        NkString codec;          // "aac","opus","vorbis","vp8","vp9","h264","pcm","mp3",...
        int32 sampleRate = 0;    int32 channels = 0;      // audio
        int32 width = 0;         int32 height = 0;        // video
        int32 bitsPerSample = 0; bool pcmBigEndian = false;
        NkVector<nk_uint8> codecPrivate;                  // OpusHead pour Opus, etc.
};
struct NkMediaInfo {
        NkMediaContainer container = NkMediaContainer::NK_UNKNOWN;
        NkVector<NkMediaTrack> tracks;
        const char* ContainerName() const;
        const NkMediaTrack* FirstAudio() const;
        const NkMediaTrack* FirstVideo() const;
};
struct NkMediaProbe {
        static bool Probe(const uint8* data, usize size, NkMediaInfo& out);
        static bool ProbeFile(const char* path, NkMediaInfo& out);
        static NkMediaContainer DetectContainer(const uint8* data, usize size);
        static bool SelfTest();
};
\end{nkcode}

L'usage réel, tiré de l'application de test :

\begin{nkcode}{Applications/NKMediaTest/src/main.cpp:331-353 (abrégé)}
    if (argc >= 2) {
        media::NkMediaInfo info;
        if (!media::NkMediaProbe::ProbeFile(argv[1], info)) {
            printf("[ERREUR] probe echoue : %s\n", argv[1]);
            return 1;
        }
        printf("Conteneur : %s\n", info.ContainerName());
        printf("Pistes    : %d\n", (int)info.tracks.Size());
        for (uint64 i = 0; i < info.tracks.Size(); ++i) {
            const media::NkMediaTrack &t = info.tracks[i];
            printf("  [%d] %-5s codec=%-6s", (int)i, TrackTypeName(t.type), t.codec.CStr());
            if (t.type == media::NkMediaTrackType::NK_AUDIO)
                printf(" %d Hz, %d canal(aux)", t.sampleRate, t.channels);
            if (t.type == media::NkMediaTrackType::NK_VIDEO)
                printf(" %dx%d", t.width, t.height);
            printf("\n");
        }
\end{nkcode}

\texttt{FirstAudio()} et \texttt{FirstVideo()} évitent de parcourir soi-même les
pistes quand on ne s'intéresse qu'à la principale — ils rendent \texttt{nullptr}
s'il n'y en a pas.

\begin{nkretenir}
Avant d'ouvrir un fichier avec \texttt{NkVideoReader}, sondez-le. Vous saurez
alors si le codec est l'un de ceux que vous savez traiter, et vous pourrez
afficher un message utile au lieu d'un échec muet. \texttt{ProbeFile} est
instantané : il ne décode aucune image.
\end{nkretenir}

\section{Écrire une vidéo}

\subsection{L'API}

\begin{nkcode}{Kernel/Runtime/NKMedia/src/NKMedia/Video/NkVideoWriter.h:26-81}
enum class NkVideoCodec { RAW_BGR, MJPEG, MPEG1 };
enum class NkVideoContainer { AVI, MOV, ELEMENTARY };
// (NkVideoInputFormat { RGB24, RGBA32, BGR24 } vit dans Video/NkVideoTypes.h:16-20)

struct NkVideoConfig {
        int32 width = 0;  int32 height = 0;
        int32 fpsNum = 30;  int32 fpsDen = 1;
        NkVideoCodec codec = NkVideoCodec::MJPEG;
        NkVideoContainer container = NkVideoContainer::AVI;
        int32 quality = 90;                 // MJPEG 1..100
        int32 audioSampleRate = 0;          // 0 = video muette
        int32 audioChannels = 0;
};
class NkVideoWriter {
        bool Open(const char* path, const NkVideoConfig& cfg);
        bool WriteFrame(const uint8* pixels, NkVideoInputFormat fmt);
        bool AddAudioSamples(const int16* interleaved, usize frameCount);
        bool Close();
        bool IsOpen() const;   int32 FrameCount() const;
};
\end{nkcode}

Notez la cadence exprimée en fraction : \texttt{fpsNum / fpsDen}. Pour du
30~images/s, \texttt{30/1} ; pour du 29,97 (NTSC), \texttt{30000/1001}. Les
formats vidéo n'aiment pas les flottants.

\subsection{Le squelette de référence}

\begin{nkcode}{Applications/NKVideoTest/src/main.cpp:83-112}
    bool MakeVideo(const char *path, media::NkVideoCodec codec, media::NkVideoContainer container, int32 w, int32 h,
                   int32 fps, int32 nframes) {
        media::NkVideoConfig cfg;
        cfg.width = w;
        cfg.height = h;
        cfg.fpsNum = fps;
        cfg.fpsDen = 1;
        cfg.codec = codec;
        cfg.container = container;
        cfg.quality = 88;

        media::NkVideoWriter vw;
        if (!vw.Open(path, cfg)) {
            ::printf("  [ERREUR] ouverture %s\n", path);
            return false;
        }
        uint8 *rgb = (uint8 *)memory::NkAlloc((usize)w * h * 3);
        for (int32 fr = 0; fr < nframes; ++fr) {
            RenderFrame(rgb, w, h, fr, nframes);
            if (!vw.WriteFrame(rgb, media::NkVideoInputFormat::RGB24)) {
                ::printf("  [ERREUR] trame %d\n", fr);
                memory::NkFree(rgb);
                vw.Close();
                return false;
            }
        }
        memory::NkFree(rgb);
        vw.Close();
        return true;
    }
\end{nkcode}

Quatre choses à retenir de ce bloc.

\begin{enumerate}
  \item \textbf{Le tampon est alloué une fois}, hors de la boucle, et réutilisé.
        Allouer une image par trame est le meilleur moyen de rendre l'encodage
        plus lent que le décodage ;
  \item \textbf{la mémoire vient de \texttt{memory::NkAlloc}} et repart par
        \texttt{memory::NkFree} — la règle du chapitre~5 ne change pas ;
  \item \textbf{le chemin d'erreur appelle quand même \texttt{Close()}} ;
  \item \textbf{le format d'entrée est explicite} :
        \texttt{NkVideoInputFormat::RGB24}. Le writer accepte aussi
        \texttt{RGBA32} et \texttt{BGR24} et fait la conversion.
\end{enumerate}

\begin{nkpiege}[\texttt{Close()} n'est pas optionnel]
Un fichier AVI se termine par un index ; un MP4 par un atome \texttt{moov}.
Aucun des deux n'est écrit tant que \texttt{Close()} n'a pas été appelé. Sans
lui, vous obtenez un fichier de la bonne taille, contenant toutes vos images, et
qu'\textbf{aucun lecteur au monde n'ouvrira}. Le message typique est
« \texttt{moov atom not found} ».

Le dépôt prend explicitement cette précaution dans son code de capture :

\begin{nkcode}{Applications/NKViewportDemo/src/NKViewportDemo/main.cpp:358-362}
    // Robustesse : si la fenêtre est fermée AVANT la fin de l'enregistrement, on finalise quand même
    // (écrit le moov) → le MP4 reste lisible (sinon "moov atom not found").
    if (rec.IsRecording()) {
        rec.End();
    }
\end{nkcode}

La même règle vaut pour \texttt{NkVideoRecorder::End()} et
\texttt{NkWebmWriter::Finalize()}. \textbf{Prévoyez une finalisation sur tous les
chemins de sortie, y compris la fermeture brutale de la fenêtre.}
\end{nkpiege}

\subsection{Deux limites de l'écriture}

\begin{itemize}
  \item \textbf{L'audio n'est pas supportée partout} :
        \nkfile{Video/NkVideoWriter.h:50-52} précise que la piste audio est
        « supportée par AVI \dots{} et MOV/MP4 \dots{} Non supportée par
        ELEMENTARY/MPEG1 (\textbf{Open échoue si demandée}) ». Au moins l'échec
        est franc ;
  \item \textbf{le sens des pixels} : les encodeurs veulent du \emph{haut en
        bas} (\nkfile{Video/NkVideoTypes.h:16-20}). Or OpenGL rend de bas en
        haut. Nous verrons le paramètre qui corrige cela.
\end{itemize}

\subsection{La séquence d'images, à la manière de Blender}

\begin{nkcode}{Kernel/Runtime/NKMedia/src/NKMedia/Video/NkImageSequenceWriter.h:21-42}
enum class NkImageSeqFormat { PNG, JPEG, BMP, TGA, QOI };
struct NkImageSequenceWriter {
        bool Open(const char* dir, const char* basename, int32 width, int32 height,
                  NkImageSeqFormat fmt, int32 padding = 4, int32 quality = 90);
        bool WriteFrame(const uint8* pixels, NkVideoInputFormat fmt);
        bool Close();  int32 FrameCount() const;
};
\end{nkcode}

\begin{nkcode}{Applications/NKVideoTest/src/main.cpp:185-201}
        media::NkImageSequenceWriter seq;
        bool ok = seq.Open(outDir, "nkframe", W, H, media::NkImageSeqFormat::PNG, 4, 90);
        if (ok) {
            uint8 *rgb = (uint8 *)memory::NkAlloc((usize)W * H * 3);
            for (int32 fr = 0; fr < 5 && ok; ++fr) { // 5 images d'exemple
                RenderFrame(rgb, W, H, fr, N);
                ok = seq.WriteFrame(rgb, media::NkVideoInputFormat::RGB24);
            }
            memory::NkFree(rgb);
            seq.Close();
        }
\end{nkcode}

Cela produit \texttt{nkframe\_0000.png}, \texttt{nkframe\_0001.png}\dots{} Cette
classe délègue entièrement aux codecs de NKImage : c'est le chemin le plus fiable
du module, et le plus facile à déboguer — vous pouvez ouvrir chaque image.

Et pour transformer ensuite la séquence en vidéo, il existe un raccourci en une
ligne :

\begin{nkcode}{Kernel/Runtime/NKMedia/src/NKMedia/Video/NkVideoConverter.h:24-30}
static int32 ImageSequenceToVideo(const char* prefix, int32 digits, const char* suffix,
                                  int32 first, int32 count, const char* outPath,
                                  int32 fpsNum, int32 fpsDen, int32 quality = 90);
static int32 MjpegAviToVideo(const char* aviPath, const char* outPath, int32 quality = 90);
\end{nkcode}

Ces deux fonctions renvoient le \emph{nombre de trames} traitées, pas un
booléen : zéro signifie l'échec.

\section{Lire une vidéo}

\subsection{L'API}

\begin{nkcode}{Kernel/Runtime/NKMedia/src/NKMedia/Video/NkVideoReader.h:26-70}
    struct NkVideoFrame {
            NkVector<nk_uint8> rgba;         // width*height*4, row-major, haut->bas
            int32 width = 0;   int32 height = 0;
            int64 timestampMs = 0;
            int32 index = -1;
    };
    struct NkVideoReaderInfo {
            int32 width = 0;   int32 height = 0;
            int32 frameCount = -1;           // -1 = inconnu
            double fps = 0.0;
            NkString codec;                  // "mjpeg"|"rawrgb"|"h264"|"hevc"|"vp8"|"vp9"|...
            NkString container;              // "avi"|"mov"|"mp4"|"sequence"
    };
    class NkVideoReader {
            NkVideoReader(const NkVideoReader&) = delete;          // NON COPIABLE
            bool Open(const char* path);
            bool IsOpen() const;
            const NkVideoReaderInfo& Info() const;
            bool ReadFrame(NkVideoFrame& out);                      // false = fini
            bool SeekFrame(int32 index);
            int32 CurrentIndex() const;
            void Close();
            static bool SelfTest();
    };
\end{nkcode}

Le point capital est dans le commentaire de \texttt{NkVideoFrame::rgba} :
\textbf{RGBA8, ligne par ligne, de haut en bas}. C'est exactement le format
qu'attendent \texttt{NkTexture::Update} et \texttt{UploadImageRGBA}. Il n'y a
aucune conversion à écrire.

\subsection{La boucle de lecture}

\begin{nkcode}{Applications/NkVideoReadTest/src/main.cpp:3041-3070 (abrégé)}
    const char *path = argv[1];
    NkVideoReader rd;
    if (!rd.Open(path)) {
        printf("  [KO] impossible d'ouvrir/lire : %s\n", path);
        return 1;
    }
    const NkVideoReaderInfo &in = rd.Info();
    printf("  conteneur=%s codec=%s  %dx%d  %.2f fps  frames=%d\n", in.container.CStr(), in.codec.CStr(), in.width,
           in.height, in.fps, in.frameCount);

    int32 count = 0;
    NkVideoFrame fr;
    while (rd.ReadFrame(fr)) {
        /* ... traiter fr.rgba ... */
        ++count;
    }
\end{nkcode}

Remarquez que \texttt{fr} est déclarée \emph{hors} de la boucle et réutilisée :
son \texttt{NkVector} conserve sa capacité d'une trame à l'autre. La déclarer
dans la boucle provoquerait une allocation par image.

\texttt{Open} accepte aussi un \textbf{dossier d'images} : le champ
\texttt{container} vaut alors \texttt{"sequence"} et le codec \texttt{"image"}.
C'est le moyen le plus simple de tester votre code de lecture sans avoir de
vidéo sous la main.

\subsection{Le seek n'est pas ce que vous croyez}

\begin{nkcode}{Applications/NkVideoReadTest/src/main.cpp:2762-2765}
 SeekFrame se contente de repositionner sur l'IDR précédente (voir implémentation) : il
 faut ensuite redécoder EN AVANT jusqu'à la cible, exactement comme le fait la boucle de
 rattrapage du lecteur (Applications/NkVideoPlayer).
\end{nkcode}

Dans un format à trames inter-dépendantes comme H.264, l'image numéro 1\,000 ne
peut pas être décodée seule : elle décrit des différences par rapport aux
précédentes. \texttt{SeekFrame(1000)} vous replace donc sur la dernière image
autonome avant la cible, et c'est à vous d'appeler \texttt{ReadFrame} jusqu'à ce
que \texttt{CurrentIndex() >= 1000}.

Sur une séquence d'images ou un AVI MJPEG — où chaque image est indépendante —
le seek est exact.

\subsection{Décoder coûte cher}

Une phrase à graver : \textbf{ne décodez jamais plus d'une image par image
affichée}. Le lecteur du dépôt borne explicitement sa boucle de rattrapage :

\begin{nkcode}{Applications/NkVideoPlayer/src/main.cpp:562-620 (extrait)}
        if (!paused && !ended) {
            if (audioOn && !streamPlayer.IsFinished())
                mediaClock = streamPlayer.GetPositionSeconds();
            else
                mediaClock += (float64)(dt * speed);
            const int32 targetIdx = (int32)(mediaClock * (float64)fps);
            /* resynchronisation dure si on a plus de 1,5 s de retard */
            const int32 kHardResyncLagFrames = (int32)(fps * 1.5f);
            if (!firstFrame && (targetIdx - reader.CurrentIndex()) > kHardResyncLagFrames) {
                if (reader.SeekFrame(targetIdx) && reader.ReadFrame(fr))
                    havePendingFrame = true;
            }
            NkChrono burstTimer;
            while ((firstFrame || reader.CurrentIndex() < targetIdx) &&
                   burstTimer.Elapsed().ToMilliseconds() < 150.0) {
                if (reader.ReadFrame(fr)) { havePendingFrame = true; firstFrame = false; }
                else if (loop) { /* SeekFrame(0) + ReadFrame */ }
                else { paused = true; break; }
            }
            if (havePendingFrame)
                pushFrame(); // upload + dessin UNE SEULE FOIS, avec la derniere image de la rafale
        }
\end{nkcode}

Trois garde-fous dans ce seul bloc : la rafale de rattrapage est limitée à
150~ms de temps mur ; au-delà d'une seconde et demie de retard on saute
carrément ; et surtout, \textbf{l'upload GPU n'a lieu qu'une fois}, avec la
dernière image de la rafale. Décoder dix images pour n'en afficher qu'une est
acceptable ; en uploader dix ne l'est pas.

\begin{nkretenir}[L'horloge maîtresse]
Notez la première ligne : quand l'audio joue, c'est \emph{lui} qui donne
l'heure (\texttt{mediaClock = streamPlayer.GetPositionSeconds()}). Ce n'est
qu'en son absence qu'on accumule le \emph{delta time}. La raison est
physiologique : l'oreille détecte un hoquet audio de 20~ms, l'œil ne voit pas
une image manquante. On synchronise donc toujours l'image sur le son, jamais
l'inverse.
\end{nkretenir}

\section{Afficher la vidéo}

\subsection{Dans une fenêtre NkCanvas}

Le principe : créer \textbf{une} texture à la taille de la vidéo, puis la mettre
à jour à chaque nouvelle image.

\begin{nkcode}{Applications/NkVideoPlayer/src/main.cpp:355-362}
    // ── 4) Texture RGBA de la taille vidéo + sprite d'affichage ───────────────
    NkTexture frameTex;
    if (!frameTex.Create(*target.GetRenderer(), (uint32)vidW, (uint32)vidH, NkColor2D{0, 0, 0, 255})) {
        logger.Error("[NkVideoPlayer] echec creation texture {0}x{1}", vidW, vidH);
        window.Close();
        return -5;
    }
    NkSprite sprite(frameTex);
\end{nkcode}

puis, à chaque image décodée :

\begin{nkcode}{Applications/NkVideoPlayer/src/main.cpp:424-431}
    auto pushFrame = [&]() {
        rawW = fr.width;
        rawH = fr.height;
        rawFrame = fr.rgba;
        applyLook(fr.rgba.Data(), (uint64)fr.width * fr.height);
        frameTex.Update(fr.rgba.Data(), (uint32)fr.width, (uint32)fr.height, 0, 0);
        haveFrame = true;
    };
\end{nkcode}

et le dessin, qui n'a rien de particulier :

\begin{nkcode}{Applications/NkVideoPlayer/src/main.cpp:664-668}
        target.Clear(NkColor2D{0, 0, 0, 255});
        if (haveFrame)
            target.Draw(static_cast<const NkDrawable &>(sprite));
        target.Display();
\end{nkcode}

\begin{nkpiege}[Ne recréez jamais la texture par image]
\texttt{Create} alloue une texture GPU ; \texttt{Update} y recopie des pixels.
Appeler \texttt{Create} à chaque trame, c'est allouer et libérer une texture
soixante fois par seconde : la mémoire GPU se fragmente et le pilote finit par
protester. Créez une fois, mettez à jour ensuite.
\end{nkpiege}

\subsection{Dans une interface}

Le motif est le même, avec l'API du backend d'interface. L'IDE du dépôt en donne
la version complète, y compris le cas du changement de dimensions :

\begin{nkcode}{Applications/NKCode/src/NKCode/Shell/NkVideoViewer.h:70-82}
        inline void NkVideoUpload(editorkit::NkEditorShell *shell, NkVideoClip *c) {
            if (!shell || c->frame.rgba.Empty() || c->frame.width <= 0 || c->frame.height <= 0)
                return;
            const uint8 *px = c->frame.rgba.Data();
            if (c->texId == 0 || c->texW != c->frame.width || c->texH != c->frame.height) {
                c->texId = shell->UploadRGBA(px, c->frame.width, c->frame.height);
                c->texW = c->frame.width;
                c->texH = c->frame.height;
            } else {
                shell->UpdateRGBA(c->texId, px, c->frame.width, c->frame.height);
            }
            c->haveFrame = c->texId != 0;
        }
\end{nkcode}

Premier appel ou changement de taille : \texttt{UploadRGBA}, qui crée. Sinon :
\texttt{UpdateRGBA}, qui recopie. C'est exactement la même logique que pour
NKImage au chapitre~5, avec une texture qui vit plus longtemps qu'une frame.

Et le cadencement, côté interface :

\begin{nkcode}{Applications/NKCode/src/NKCode/Shell/NkVideoViewer.h:153-170}
            float32 dt = ctx.input.dt;
            if (dt <= 0.f || dt > 0.25f)
                dt = 1.f / 60.f; // garde-fou (onglet revenu au premier plan / hoquet)
            const float32 frameDur = 1.f / (c->fps * (c->speed > 0.01f ? c->speed : 0.01f));
            if (c->playing && !c->ended) {
                c->acc += dt;
                if (c->acc > frameDur * 4.f)
                    c->acc = 0.f; // evite le rattrapage explosif
                while (c->acc >= frameDur) {
                    c->acc -= frameDur;
                    if (c->reader->ReadFrame(c->frame)) {
                        c->index = c->frame.index;
                        NkVideoUpload(shell, c);
                    }
                    /* ... loop / fin ... */
                }
            }
\end{nkcode}

Les deux garde-fous méritent d'être notés, parce qu'ils correspondent à deux
situations réelles : un \emph{delta time} aberrant quand l'onglet revient au
premier plan après plusieurs secondes, et un accumulateur qui a débordé et
voudrait décoder quarante images d'un coup.

\section{Enregistrer ce que l'application affiche}

\subsection{L'API}

\begin{nkcode}{Kernel/Runtime/NKMedia/src/NKMedia/Video/NkVideoRecorder.h:39-77}
enum class NkRecorderCodec { H264, MJPEG };

struct NkVideoRecorder {
        bool Begin(const char* path, int32 width, int32 height, int32 fpsNum = 60, int32 fpsDen = 1,
                   int32 qp = 20, int32 maxQueuedFrames = 32,
                   NkRecorderCodec codec = NkRecorderCodec::H264, int32 mjpegQuality = 90);
        int32 AddAudio(int32 sampleRate, int32 channels, const char* lang3 = nullptr);
        int32 AddSubtitleTrack(const char* lang3 = nullptr);
        bool PushVideo(const uint8* pixels, NkVideoInputFormat fmt, bool flipVertical = false);
        void PushAudio(int32 trackIdx, const int16* interleaved, uint32 frames);
        void AddSubtitle(int32 trackIdx, const char* utf8, uint32 startMs, uint32 durMs);
        bool End();
        bool IsRecording() const;  int32 FrameCount() const;
        int32 QueueDepth();  uint64 DroppedFrames();  double EncodeFps();
};
\end{nkcode}

La différence avec \texttt{NkVideoWriter} est architecturale :

\begin{nkcode}{Kernel/Runtime/NKMedia/src/NKMedia/Video/NkVideoRecorder.h:5-8}
 **Encodage sur un THREAD de fond** : la boucle de rendu pousse les trames capturées
 dans une file (rapide) et un worker encode en arrière-plan → l'application reste
 FLUIDE pendant la capture (l'encodeur H.264 est lourd).
\end{nkcode}

\texttt{PushVideo} rend la main immédiatement ; l'encodage a lieu ailleurs.
\texttt{End()} joint ce fil : \textbf{ne détruisez jamais un recorder sans avoir
appelé \texttt{End()}}.

\subsection{La capture, de bout en bout}

Le point de jonction avec le reste du moteur est délicieusement simple : la
capture d'écran est une \texttt{NkImage}.

\begin{nkcode}{Applications/NKViewportDemo/src/NKViewportDemo/main.cpp:329-362 (abrégé)}
        if (record && target->CaptureToImage(capImg) && capImg.IsValid()) {
            if (!rec.IsRecording()) {
                if (rec.Begin("viewport_capture.mp4", capImg.Width(), capImg.Height(), 30, 1, 24)) {
                    recAudio = rec.AddAudio(kRate, 1, "fre");
                    const int32 st = rec.AddSubtitleTrack("fre");
                    rec.AddSubtitle(st, "Capture live NKCanvas (DX11) - Nkentseu", 0, 4000);
                }
            }
            if (rec.IsRecording()) {
                rec.PushVideo(capImg.Pixels(), media::NkVideoInputFormat::RGBA32);
                rec.PushAudio(recAudio, recTone.Data(), (uint32)kAudioPerFrame);
                if (recFrames >= kRecTotal) { rec.End(); running = false; }
            }
        }
\end{nkcode}

\texttt{NkRenderWindow::CaptureToImage(NkImage\&)} remplit une \texttt{NkImage},
et \texttt{capImg.Pixels()} part directement dans \texttt{PushVideo}. NKImage,
NkCanvas et NKMedia se rejoignent en une ligne.

\subsection{Trois pièges du recorder}

\begin{nkpiege}[L'ordre des appels]
\nkfile{Video/NkVideoRecorder.h:51-54} : \texttt{AddAudio} et
\texttt{AddSubtitleTrack} sont à appeler « juste après \texttt{Begin} (avant les
trames) ». Ajouter une piste en cours d'enregistrement n'a pas de sens dans un
conteneur MP4, dont la table des pistes est écrite en tête.
\end{nkpiege}

\begin{nkpiege}[En MJPEG, l'audio est ignorée — silencieusement]
\begin{nkcode}{Kernel/Runtime/NKMedia/src/NKMedia/Video/NkVideoRecorder.h:36-38}
 conteneur MOV/MP4, **VIDEO SEULE (audio/sous-titres ignores dans ce mode)**
\end{nkcode}

C'est le piège le plus vicieux du module : \texttt{AddAudio} vous rend un index
parfaitement valide, \texttt{PushAudio} accepte vos échantillons sans broncher,
et le fichier produit est muet. Rien ne vous prévient.

Vous êtes donc devant un choix : MJPEG (léger, sûr, muet) ou H.264 (avec audio,
mais l'encodeur est marqué « en cours »). Pour apprendre, prenez MJPEG et
enregistrez l'audio à part.
\end{nkpiege}

\begin{nkpiege}[Le recorder abandonne des trames]
\begin{nkcode}{Kernel/Runtime/NKMedia/src/NKMedia/Video/NkVideoRecorder.h:45-46}
 `maxQueuedFrames` BORNE la file video : si l'encodeur (H.264 lourd) prend du retard,
 les nouvelles trames sont ABANDONNEES (drop-newest) au lieu de gonfler la memoire
 (temps reel : perdre une trame vaut mieux que geler). Abandons comptes (DroppedFrames()).
\end{nkcode}

C'est le bon comportement — mais il faut le savoir, sinon on s'étonne qu'une
capture de dix secondes à 60~images/s en contienne 400. Surveillez
\texttt{DroppedFrames()} et \texttt{QueueDepth()}, et régulez si nécessaire :

\begin{nkcode}{Applications/Sandbox/src/Demo/main.cpp:1025-1050 (extrait)}
                    const bool encoderBusy = recorder && recorder->QueueDepth() >= 24;
                    if (!encoderBusy)
                        (void)recordCapture.EnqueueCopy(...);
\end{nkcode}
\end{nkpiege}

\subsection{Le sens des pixels}

\begin{nkcode}{Kernel/Runtime/NKMedia/src/NKMedia/Video/NkVideoRecorder.h:56-57}
 `flipVertical` pour framebuffer bottom-up (OpenGL)
\end{nkcode}

OpenGL considère l'origine en bas à gauche ; les formats vidéo la placent en
haut à gauche. Si votre capture provient d'un \emph{framebuffer} OpenGL lu
directement, passez \texttt{flipVertical = true}. Le symptôme d'un oubli est
sans ambiguïté : la vidéo est à l'envers.

\section{Le démux audio, en deux mots}

Nous n'en aurons pas besoin dans ce cours, mais vous croiserez la classe :

\begin{nkcode}{Kernel/Runtime/NKMedia/src/NKMedia/NkMediaDemux.h:22-48}
struct NkMediaPacket {
        usize offset = 0;   // position dans le buffer SOURCE (ne copie pas !)
        usize size = 0;
        int64 timestampMs = 0;
        int64 granule = -1;             // OGG
        int64 discardPaddingNs = 0;     // WebM
};
struct NkMediaDemux {
        static bool ExtractAudioPackets(const uint8* data, usize size, const NkMediaInfo& info,
                                        NkVector<NkMediaPacket>& out);
        static bool ExtractAudioPacketsFile(const char* path, NkVector<nk_uint8>& outBytes,
                                            NkMediaInfo& outInfo, NkVector<NkMediaPacket>& out);
        static bool SelfTest();
};
\end{nkcode}

\begin{nkpiege}[Les paquets pointent dans le tampon source]
\texttt{NkMediaPacket} ne contient pas de données : il contient un
\emph{décalage}. Le commentaire est explicite — « Un paquet encodé (pointe dans
le buffer d'origine ; ne copie pas) » (\nkfile{NkMediaDemux.h:22}), et pour la
variante fichier : « \texttt{outBytes} reçoit le buffer (\textbf{à garder
vivant} car les paquets pointent dedans) » (\nkfile{NkMediaDemux.h:47}).

Si votre \texttt{NkVector<nk\_uint8> bytes} sort de portée avant que vous ayez
fini de lire les paquets, vous lisez de la mémoire libérée. Gardez les deux
ensemble, dans la même portée.
\end{nkpiege}

\section{Vérifier que tout marche sur votre machine}

NKMedia est le module le mieux doté du moteur en auto-tests — il en compte 52.
L'application \texttt{NkVideoReadTest}, lancée \emph{sans argument}, les
enchaîne :

\begin{nkcode}{Applications/NkVideoReadTest/src/main.cpp:252-276}
    if (argc < 2) {
        printf("  [self-test] ecrire AVI MJPEG -> relire -> verifier...\n");
        bool ok = NkVideoReader::SelfTest();
        printf("  [ %s ] NkVideoReader::SelfTest (AVI MJPEG round-trip)\n", ok ? "OK " : "KO");
        bool okH264 = NkH264Decoder::SelfTest();
        bool okCavlc = NkH264Cavlc::SelfTest();
        bool okVp9 = NkVp9Decoder::SelfTest();
        bool okAv1 = NkAv1Decoder::SelfTest();
        bool okHevc = NkHevcDecoder::SelfTest();
        bool okHevcCabac = NkHevcCabacState::SelfTest();
        bool okWav = NkWavWriter::SelfTest();
        bool okWebm = NkWebmWriter::SelfTest();
        bool all = ok && okH264 && okCavlc && okVp9 && okAv1 && okHevc && okHevcCabac && okWav && okWebm;
        printf("=== %s ===\n", all ? "LECTURE VIDEO OPERATIONNELLE" : "ECHEC");
        return all ? 0 : 1;
    }
\end{nkcode}

Le premier de la liste est le plus parlant : il écrit un AVI MJPEG, le relit, et
vérifie que ce qui sort correspond à ce qui est entré. C'est ce
\emph{round-trip} qui valide indirectement \texttt{NkVideoWriter}, qui n'a pas
d'auto-test propre — pas plus que \texttt{NkVideoRecorder},
\texttt{NkVideoConverter}, \texttt{NkImageSequenceWriter}, ni les écrivains de
conteneurs.

Deux détails pratiques :

\begin{itemize}
  \item ces auto-tests \textbf{écrivent dans le répertoire courant} :
        \texttt{nkvideoreader\_selftest.avi},
        \texttt{nkwavwriter\_selftest.wav},
        \texttt{nkwebmwriter\_selftest.webm}. Lancez-les depuis un dossier de
        travail, pas depuis la racine du dépôt ;
  \item le décodeur HEVC lit une variable d'environnement,
        \texttt{NK\_HEVC\_THREADS} : absente ou à 0, il choisit tout seul ; à 1,
        il travaille séquentiellement (\nkfile{NkVideoReadTest/src/main.cpp:249-250}).
        Utile pour départager un bug de décodage d'un bug de parallélisme.
\end{itemize}

\section{Une leçon de performance qui dépasse la vidéo}

La feuille de route documente un bug de performance instructif :

\begin{nkcode}{Kernel/Runtime/NKMedia/ROADMAP.md:1447-1458 (abrégé)}
 le vrai coupable etait `NkVideoReader::ReadFrame` [...] : la gestion du DPB (liste de
 references) COPIAIT EN PROFONDEUR (au lieu de deplacer) jusqu'a 16 `NkH264Frame`
 (plans Y/Cb/Cr + grilles de mouvement, ~380 Ko/image) DEUX FOIS par frame decodee
 [...] cout mesure ~80 ms/frame et CROISSANT [...] >90% du temps total. Fix : traits::NkMove
\end{nkcode}

Quatre-vingts millisecondes par image, dont plus de 90\,\% en copies inutiles.
La leçon est transférable bien au-delà de la vidéo : \textbf{dans une boucle
chaude, une copie profonde qui aurait dû être un déplacement peut représenter
l'essentiel du temps d'exécution}. Le champ \texttt{NkVideoFrame::rgba} est un
conteneur : le copier par valeur à chaque image coûte cher, et le lecteur du
dépôt ne le fait que parce qu'il a explicitement besoin d'une copie non
retouchée (\nkfile{Applications/NkVideoPlayer/src/main.cpp:427}).

Le même document ajoute un avertissement de méthode, qui vaut pour toute mesure :

\begin{nkcode}{Kernel/Runtime/NKMedia/ROADMAP.md:1449-1452}
 (mesure fiable : chrono mur autour d'un run borne, PAS le champ `t=` de `NkVideoReadTest`
 qui affiche le PTS video, pas le temps de decodage reel — piege rencontre une fois)
\end{nkcode}

Le \emph{PTS} est l'instant auquel une image doit être affichée dans la vidéo. Il
n'a rien à voir avec le temps qu'il a fallu pour la décoder. Confondre les deux,
c'est mesurer la durée du film au lieu de la vitesse du lecteur.

\begin{nkbilan}
NKMedia lit et écrit des vidéos, entièrement depuis zéro. Avant d'ouvrir un
fichier, on le \textbf{sonde} : ce qu'il contient réellement ne se déduit pas de
son extension. À l'affichage, une vidéo n'est qu'une texture qu'on remplace à
chaque image --- le rendu n'a aucune notion de « vidéo ». Le \texttt{Close()}
n'est pas une politesse : sans lui le fichier écrit est \emph{incomplet}, et il
s'ouvre parfois quand même, ce qui rend l'oubli difficile à détecter.
\end{nkbilan}

\section{Exercices}

\begin{nkdemo}[1 — Sonder avant d'ouvrir]
Écrivez un petit outil en ligne de commande qui prend un chemin, appelle
\texttt{NkMediaProbe::ProbeFile}, et affiche le conteneur et toutes les pistes.
Passez-lui successivement un WAV, un MP3, un MP4 et un fichier texte renommé en
\texttt{.mp4}.

Ajoutez ensuite une règle : si la première piste vidéo n'a pas un codec parmi
\texttt{"mjpeg"}, \texttt{"rawrgb"} ou \texttt{"image"}, affichez un
avertissement disant que la lecture n'est pas garantie par ce cours. Vous venez
d'écrire le garde-fou qui manque à la plupart des applications.
\end{nkdemo}

\begin{nkdemo}[2 — Le \emph{round-trip} complet]
Générez une vidéo de 100 images en MJPEG/AVI, avec un motif dont vous pouvez
prédire le contenu — par exemple un rectangle qui se déplace d'un pixel par
image. Relisez-la ensuite avec \texttt{NkVideoReader} et vérifiez que le nombre
d'images correspond, et que le rectangle est bien là où vous l'attendez à
l'image 50.

Recommencez avec le conteneur MOV. Puis essayez \texttt{RAW\_BGR} et comparez la
taille des fichiers : vous aurez une idée concrète de ce que compresser veut
dire.
\end{nkdemo}

\begin{nkdemo}[3 — Prouver que \texttt{Close()} est obligatoire]
Reprenez l'exercice précédent et commentez l'appel à \texttt{vw.Close()}.
Regardez la taille du fichier produit — elle est presque identique — puis
essayez de l'ouvrir avec votre lecteur vidéo habituel, et avec
\texttt{NkVideoReader}.

Cet exercice prend deux minutes et vous fera gagner une soirée le jour où votre
application se fermera brutalement en cours d'enregistrement.
\end{nkdemo}

\begin{nkpratique}[4 — Un lecteur cadencé, dans votre interface]
Ajoutez à votre application NKGui un panneau qui lit une séquence d'images
depuis un dossier (c'est le chemin le plus simple : \texttt{Open} accepte un
dossier) et l'affiche avec le widget \texttt{Image} du chapitre~5.

Cadencez sur \texttt{Info().fps} avec un accumulateur, en reprenant les deux
garde-fous de \texttt{NkVideoViewer} : \emph{delta time} aberrant et
accumulateur débordé. Ajoutez un bouton pause et un \texttt{SliderFloat} de
vitesse. Vérifiez que vous n'appelez \texttt{UploadRGBA} qu'une seule fois — au
premier affichage — et \texttt{UpdateRGBA} ensuite.
\end{nkpratique}

\begin{nkquestions}
\begin{enumerate}
\item Pourquoi sonder un fichier plutôt que se fier à son extension ?
\item Comment une image décodée arrive-t-elle à l'écran ?
\item Que manque-t-il à un fichier vidéo dont on a oublié le \texttt{Close()} ?
\item Où se branche un traitement des pixels dans la chaîne de lecture ?
\item Qu'est-ce qui cadence la lecture, et que se passe-t-il si on l'ignore ?
\end{enumerate}
\end{nkquestions}

\begin{nklogique}
Une vidéo que vous venez d'écrire s'ouvre dans votre lecteur, mais pas dans celui
du système.

\begin{enumerate}
\item Est-ce forcément un défaut de votre écriture ? Donnez au moins une cause
      où ce n'en est pas un.
\item Que peut faire votre lecteur qu'un lecteur du système ne fait pas, et
      pourquoi cette indulgence est-elle un piège pour vous ?
\item Concluez sur la règle générale : pourquoi un producteur ne doit-il
      \emph{jamais} être validé par son propre consommateur seul ?
\end{enumerate}
\end{nklogique}
